The feature extractor for MNIST image data is given in Table~\ref{demand-mnist-feature}, which is used for both stage 1 and 2. For DeepIV, we used the original structure proposed in \citet{Hartford2017}, which is described in Table~\ref{deepiv-demand-network-mnist}. We follow the default dropout rate in \citet{Hartford2017}, which depends on the data size. For DFIV, we used the structure described in Table~\ref{dfiv-demand-network-mnist}. The regularizer $\lambda_1, \lambda_2$ are both set to 0.1 as a result of the tuning procedure described in Appendix~\ref{sec:hyper-param}. For KIV, we used the Gaussian kernel where the bandwidth is determined by the median trick and also used random Fourier features with 100-dimensions. For DeepGMM, we used the same structure as DeepIV but no dropout is applied and the last layer of instrument net is changed to a fully connected layer which maps 32 dimensions to 1 dimension.

\begingroup
\renewcommand{\arraystretch}{1.2}
\begin{table}[t]
    \caption{Network structures of feature extractor used in demand design experiment with MNIST. For each convolution
    layer, we list the input dimension, output dimension, kernel size, stride, and padding. For the input layer, we provide the input variable. For the fully-connected layers (FC), we provide the input and output dimensions. For max-pool, we list the size of the kernel. Dropout rate here is set to 0.1. SN denotes Spectral Normalization \citep{Miyato2018}.}
    
    \label{demand-mnist-feature}
    \centering
    \begin{minipage}{0.8\hsize}
        \centering
        \begin{tabular}{cc}
            \multicolumn{2}{c}{\bf ImageFeature} 
            \\ \hline
            Layer & Configuration \\ \hline
            1 & Input($S$) \\ \hline
            2 & Conv2D (1, 64, 3, 1, 1), ReLU, SN \\ \hline
            3 & Conv2D (64, 64, 3, 1, 1), ReLU, SN \\ \hline
            4 & MaxPool(2, 2) \\ \hline
            5 & Dropout \\ \hline
            6 & FC(9216, 64), ReLU \\ \hline
            7 & Dropout \\ \hline
            8 & FC(64, 32), ReLU \\ \hline
            \end{tabular}    
    \end{minipage}
\end{table}
\endgroup  


\begingroup
\renewcommand{\arraystretch}{1.2}
\begin{table}[t]
    \caption{Network structures of DeepIV in demand design with MNIST data.  For the input layer, we provide the input variable. For the fully-connected layers (FC), we provide the input and output dimensions. For mixure Gaussian output, we report the number of components. ImageFeature denotes the module given in Table~\ref{demand-mnist-feature}. Dropout rate is described in the main text. }
    
    \label{deepiv-demand-network-mnist}
    \centering
    
    \begin{minipage}{0.48\hsize}
        \centering
        \begin{tabular}{cc}
            \multicolumn{2}{c}{\bf Instrument Net} 
            \\ \hline
            Layer & Configuration \\ \hline
            1 & Input($C, T, \mathtt{Image Feature}(S)$) \\ \hline
            2 & FC(66, 32), ReLU \\ \hline
            3 & Dropout \\ \hline
            4 & MixtureGaussian(10) \\ \hline
            \end{tabular}    
    \end{minipage}
    \hfill
    \begin{minipage}{0.48\hsize}
        \centering
        \begin{tabular}{cc}
            \multicolumn{2}{c}{\bf Treatment Net} 
            \\ \hline
            Layer & Configuration \\ \hline
            1 & Input($P, T, \mathtt{Image Feature}(S)$) \\ \hline
            2 & FC(66, 32), ReLU \\ \hline
            3 & Dropout \\ \hline
            4 & FC(32, 1), ReLU \\ \hline
            \end{tabular}    
    \end{minipage}
    
\end{table}
\endgroup  


\begingroup
\renewcommand{\arraystretch}{1.2}
\begin{table}[t]
    \caption{Network structures of DFIV in demand design with MNIST. For the input layer, we provide the input variable. For the fully-connected layers (FC), we provide the input and output dimensions.  ImageFeature denotes the module given in Table~\ref{demand-mnist-feature}.}
    
    \label{dfiv-demand-network-mnist}
    
    \begin{minipage}{0.48\hsize}
        \centering
        \begin{tabular}{cc}
            \multicolumn{2}{c}{\textbf{Instrumental Feature} $\vec{\phi}_{\theta_Z}$ }
            \\ \hline
            Layer & Configuration \\ \hline
            1 & Input($C, T, \mathtt{ImageFeature}(S)$) \\ \hline
            2 & FC(66, 32), BN, ReLU\\ \hline
            \end{tabular}    
    \end{minipage}
    \hfill
    \begin{minipage}{0.48\hsize}
        \centering
        \begin{tabular}{cc}
            \multicolumn{2}{c}{\textbf{Treatment Feature} $\vec{\psi}_{\theta_X}$} 
            \\ \hline
            Layer & Configuration \\ \hline
            1 & Input($P$) \\ \hline
            2 & FC(1, 16), ReLU \\ \hline
            3 & FC(16, 1) \\ \hline
            \end{tabular}    
    \end{minipage}\\
    \centering
    
    \begin{minipage}{0.48\hsize}
        \centering
        \begin{tabular}{cc}
            \multicolumn{2}{c}{\textbf{Obsevable Feature Net} $\vec{\xi}_{\theta_O}$} 
            \\ \hline
            Layer & Configuration \\ \hline
            1 & Input($T, \mathtt{ImageFeature}(S)$) \\ \hline
            2 & FC(65, 32), BN, ReLU \\ \hline
            \end{tabular}    
    \end{minipage}
\end{table}
\endgroup  

